\section{Height-one Stieltjes coefficients and the inherited Gamma bridge}
\label{sec:gamma-foundation}

The sixth and seventh reductions use a height-one Laurent coefficient.
We record its analytic definition and prove the inherited integral bridge
in the exact convention needed here. These coefficient formulas are
established in the incoming fifth-moment report and are related to the
height-one series of Young \cite{CanonicalFifth,Young2024}; their inclusion
does not constitute a new evaluation of generalized Stieltjes constants.

For depth $d\ge1$, let
\[
 Z_d(s)=\zeta(s,\underbrace{1,\ldots,1}_{d-1})
       =\sum_{n_1>\cdots>n_d\ge1}
          \frac{1}{n_1^s n_2\cdots n_d},\qquad \re s>1.
\]
Define $c_j$ and $\lambda_{d,j}$ by
\begin{align}
 \frac1{\Gamma(1+u)}&=\sum_{j\ge0}c_ju^j,\label{eq:gamma-c-foundation}\\
 Z_d(1+u)&=\sum_{j=0}^d c_ju^{j-d}
                    +\sum_{j\ge1}\lambda_{d,j}u^j,
 \qquad \lambda_d:=\lambda_{d,1}.
 \label{eq:height-one-foundation}
\end{align}
The second expansion will follow from the next proposition. In particular,
$\lambda_{1,j}=(-1)^j\gamma_j/j!$ and $\lambda_1=-\gamma_1$.
The multiple coefficients $\lambda_{d,j}$ have a specified depth; they
should not be identified with ordinary Hurwitz coefficients $\gamma_j(a)$
without a further identity.

For $v>0$, set
\[
 B(v)=-\log(1-e^{-v}),\qquad
 V_d(v)=B(v)^d-\boldsymbol1_{(0,1)}(v)(-\log v)^d,
\]
and put
\[
 I_{d,r}=\int_0^\infty\frac{(\log v)^rV_d(v)}v\,dv
 \quad(r\ge0).
\]
At zero, $V_d(v)=O(v(1+|\log v|^{d-1}))$; at infinity it
decays exponentially. All these integrals converge absolutely.

\begin{proposition}[Fixed Mellin normalization]
\label{prop:height-one-foundation}
For $d\ge1$, as a meromorphic identity in $\re u>-1$,
\begin{equation}\label{eq:height-one-Mellin-foundation}
 \Gamma(1+u)Z_d(1+u)
 =u^{-d}+\frac{u}{d!}
          \int_0^\infty v^{u-1}V_d(v)\,dv.
\end{equation}
Consequently the coefficients in \eqref{eq:height-one-foundation} satisfy
\begin{equation}\label{eq:height-one-allcoeff-foundation}
 \lambda_{d,j}=c_{d+j}
       +\frac1{d!}\sum_{r=0}^{j-1}
                     \frac{c_{j-1-r}}{r!}I_{d,r},\qquad j\ge1.
\end{equation}
If $\Lambda_d(u)=\Gamma(1+u)Z_d(1+u)-u^{-d}$, then
$\Lambda_d^{(j)}(0)=jI_{d,j-1}/d!$ for $j\ge1$.
\end{proposition}

\begin{proof}
The strict-index generating series with $d-1$ unit indices is
$(-\log(1-y))^{d-1}/(d-1)!$. Termwise Mellin integration for
$\re u>0$ therefore gives
\[
 \Gamma(1+u)Z_d(1+u)
  =\frac1{(d-1)!}\int_0^\infty
            \frac{v^u B(v)^{d-1}}{e^v-1}\,dv
  =\frac{u}{d!}\int_0^\infty v^{u-1}B(v)^d\,dv.
\]
For the second equality integrate by parts, using
$B'(v)=-(e^v-1)^{-1}$; both boundary terms vanish in this half-plane.
Subtract the exact integral
$\int_0^1v^{u-1}(-\log v)^d\,dv=d!/u^{d+1}$.
The estimates above give holomorphy of the remaining integral for
$\re u>-1$, locally uniformly with every spectral derivative.
Expansion at zero and multiplication by \eqref{eq:gamma-c-foundation}
prove all the coefficient identities.
\end{proof}

This calculation makes the Stieltjes content of the later formulas
explicit: the spectral derivative data are ordinary convergent logarithmic
Mellin integrals, with the subtraction at $v=1$ fixed in advance.

Define the Gamma family
\begin{equation}\label{eq:Cgamma-foundation}
 G(x,t)=\frac{\Gamma(1-t)\Gamma(x)}{\Gamma(x-t)}
       =\sum_{m\ge0}C_m(x)\frac{t^m}{m!}.
\end{equation}
At positive integer $x=n$, it is the familiar complete harmonic-star
generator $\prod_{k=1}^{n-1}(1-t/k)^{-1}$; this classical connection
is also discussed by Sesma \cite{Sesma}. At $x=0$, the residue of
$G(x,t)$ is $-t$. Thus every $C_m$ with $m\ge2$ has a removable
singularity there. Expansion of the constant term gives
$C_m(0)=-m!\zeta(m)$, while $G(1,t)=1$ gives $C_m(1)=0$.

\begin{proposition}[Inherited ordinary Gamma integral]
\label{prop:Cgamma-integral-foundation}
For every $m\ge2$,
\begin{equation}\label{eq:Cgamma-integral-foundation}
 \int_0^1C_m(x)\,dx=m!(c_m-\lambda_{m-1}).
\end{equation}
The integral is ordinary and the complete integrand is continuous on
$[0,1]$.
\end{proposition}

\begin{proof}
Use $F(x,t)=G(x,-t)$ and write $h_1(x)=\psi(x)+\gamma$. The function
\[
 \Psi(t)=\int_0^1\{F(x,t)-1+t h_1(x)\}\,dx
\]
is holomorphic near zero. Indeed the residues $t/x$ and $-t/x$
cancel, leaving a jointly regular expression for $x\in[0,1]$ and
$t$ in a sufficiently small closed disc. Hence
\[
 \Psi(t)=\sum_{m\ge2}\frac{(-1)^m t^m}{m!}
                                  \int_0^1C_m(x)\,dx.
\]
For $\re t>0$, Euler's beta integral for $F$ and the ordinary
digamma integral, combined before integration, yield
\[
 \Psi(t)=t\int_0^\infty
   \left\{\frac{(1-e^{-v})^t-1}{v}
                    +\frac{e^{-v}}{1-e^{-v}}\right\}\,dv-1.
\]
Fubini is justified for the combined integrand: at zero it is
$O(v^{\re t-1}+1)$ after cancellation, and at infinity integration
in $x$ leaves an exponentially decaying numerator divided by $v$.
To put this in the fixed subtraction convention, integrate from a
lower cutoff $\delta$, use
$\int_\delta^\infty(e^v-1)^{-1}\,dv=-\log(1-e^{-\delta})$,
and compare with $\int_\delta^1(v^t-1)\,dv/v$. Taking
$\delta\downarrow0$ gives
\begin{equation}\label{eq:gamma-Psi-foundation}
 \Psi(t)=t\int_0^\infty
 \frac{(1-e^{-v})^t-1
       -\boldsymbol1_{(0,1)}(v)(v^t-1)}v\,dv.
\end{equation}
Near zero, $(1-e^{-v})^t-v^t=O(v^{\re t+1})$ uniformly on compact
parameter sets; at infinity the numerator decays exponentially.
Thus this equality continues holomorphically to a neighborhood of zero,
and coefficients may be extracted under the integral. Its degree-$m$
coefficient is $(-1)^{m-1}I_{m-1,0}/(m-1)!$.
Comparing the two expressions for $\Psi$ gives
$\int_0^1C_m=-mI_{m-1,0}$. Proposition~\ref{prop:height-one-foundation}
then proves \eqref{eq:Cgamma-integral-foundation}.
\end{proof}

Individual terms in a differential expansion of $C_m$ can have divergent
improper integrals even though $C_m$ itself is regular. The finite-part
normal form in the next section is what permits those terms to be
separated with their exact endpoint contributions retained.
